% Intended LaTeX compiler: pdflatex
\documentclass[a4paper,12pt]{article}
\usepackage[utf8]{inputenc}
\usepackage[T1]{fontenc}
\usepackage{graphicx}
\usepackage{longtable}
\usepackage{wrapfig}
\usepackage{rotating}
\usepackage[normalem]{ulem}
\usepackage{amsmath}
\usepackage{amssymb}
\usepackage{capt-of}
\usepackage{hyperref}
\usepackage[margin=0.2in]{geometry}
\usepackage{amsmath}
\usepackage{amsfonts}
\usepackage{indentfirst}
\usepackage{pgffor}
\usepackage{amssymb}
\usepackage{cancel}
\usepackage{amsthm}
\usepackage{polynom}
\usepackage{tikz}
\usepackage[tikz]{bclogo}
\usepackage{listings}
\usepackage[most]{tcolorbox}
\usepackage{forest}
\usepackage{adjustbox}
\usepackage{tikz-3dplot}
\usepackage{mathtools}
\usepackage{pgfplots}
\usetikzlibrary{tikzmark,calc,fit,matrix,arrows,automata,positioning}
\usepackage{centernot}
\usepackage{lmodern}
\usepackage{physics}
\usepackage[boxed,linesnumbered,vlined]{algorithm2e}
\usepackage{algpseudocode}
\usepackage{algorithmicx}
\usepackage{svg}
\tikzstyle{every picture}+=[remember picture]
\DeclareMathOperator{\spn}{span}
\DeclareMathOperator{\dom}{domain}
\DeclareMathOperator{\ran}{range}
\DeclareMathOperator{\rng}{range}
\DeclareMathOperator{\img}{Im}
\DeclareMathOperator{\adj}{adj}
\newcommand\dif[1]{\:\textrm{d}#1}
\DeclarePairedDelimiter\ceil{\lceil}{\rceil}
\DeclarePairedDelimiter\floor{\lfloor}{\rfloor}
\newcommand{\ihat}{\hat{\textbf{\i}}}
\newcommand{\jhat}{\hat{\textbf{\j}}}
\newcommand{\khat}{\hat{\textbf{k}}}
\newcommand{\rhat}{\hat{\textbf{r}}}
\newcommand{\thetahat}{\boldsymbol{\hat{\theta}}}
\author{mahmood}
\date{\today}
\title{digital systems homework 6}
\hypersetup{
 pdfauthor={mahmood},
 pdftitle={digital systems homework 6},
 pdfkeywords={},
 pdfsubject={},
 pdfcreator={Emacs 30.0.50 (Org mode 9.6)}, 
 pdflang={English}}
\begin{document}

\maketitle
\definecolor{Brown}{cmyk}{0,0.81,1,0.60}
\definecolor{OliveGreen}{cmyk}{0.64,0,0.95,0.40}
\definecolor{CadetBlue}{cmyk}{0.62,0.57,0.23,0}
\definecolor{lightlightgray}{gray}{0.9}
\lstset{
  language=C,                             % Code langugage
  basicstyle=\ttfamily,                   % Code font, Examples: \footnotesize, \ttfamily
  keywordstyle=\color{OliveGreen},        % Keywords font ('*' = uppercase)
  commentstyle=\color{gray},              % Comments font
  numbers=left,                           % Line nums position
  numberstyle=\tiny,                      % Line-numbers fonts
  stepnumber=1,                           % Step between two line-numbers
  numbersep=5pt,                          % How far are line-numbers from code
  backgroundcolor=\color{lightlightgray}, % Choose background color
  frame=single,                           % A frame around the code
  tabsize=2,                              % Default tab size
  captionpos=b,                           % Caption-position = bottom
  breaklines=true,                        % Automatic line breaking?
  breakatwhitespace=false,                % Automatic breaks only at whitespace?
  showspaces=false,                       % Dont make spaces visible
  showtabs=false,                         % Dont make tabls visible
  columns=flexible,                       % Column format
  morekeywords={__global__, __device__},  % CUDA specific keywords
  showstringspaces=false,                 % dont show spaces in strings
}
% so that latex doesnt complain about sage not being a language
\lstdefinelanguage{sage}{}

% environment definition for special blocks
\theoremstyle{definition}
\newtheorem{my_example}{Example}
\counterwithin{my_example}{section}
\counterwithin{my_example}{subsection}
\newtheorem{definition}{Definition}
\counterwithin{definition}{section}
\counterwithin{definition}{subsection}
\newtheorem{characteristic}{Characteristic}
\newtheorem{entailment}{Entailment}
%\newtheore*{proof}{Proof}
\newtheorem{lemma}{Lemma}
\newtheorem{question}{Question}
\newtheorem{subquestion}{Subquestion}[question]
%\newtheorem{note}{Note}
\newtheorem{answer}{Answer}
\counterwithin{answer}{question}
\counterwithin{answer}{subquestion}
\newtheorem{step}{Step}[answer]

% the following snippet auto resizes images to fit properly
% Determine if the image is too wide for the page.
% \makeatletter
% \def\ScaleIfNeeded{%
%   \ifdim\Gin@nat@width>\linewidth
%     \linewidth
%   \else
%     \Gin@nat@width
%   \fi
% }
% \makeatother
% % Resize figures that are too wide for the page.
% \let\oldincludegraphics\includegraphics
% \renewcommand\includegraphics[2][]{%
%   \oldincludegraphics[width=\ScaleIfNeeded]{#2}
% }

\newenvironment{note}
  {\begin{bclogo}[logo=\bcinfo, couleurBarre=orange, couleur=lightgray]{ note}}
  {\end{bclogo}}
\newenvironment{attention}
  {\begin{bclogo}[logo=\bcattention, couleurBarre=red, noborder=true, 
                couleur=red!15]{Important!}}
  {\end{bclogo}}

\begin{question}
given the function \(f(w,x,y,z) = \sum(5,7,8,9,10,15)\)\\[0pt]
\begin{subquestion}
write a \href{20230116122956-truth_table.org}{truth table} for the function\\[0pt]
\begin{answer}
\begin{center}
\begin{tabular}{rrrrrrll}
 & w & x & y & z & f & maxterm & minterm\\[0pt]
\hline
0 & 0 & 0 & 0 & 0 & 0 &  & w'x'y'z'\\[0pt]
1 & 0 & 0 & 0 & 1 & 0 &  & w'x'y'z\\[0pt]
2 & 0 & 0 & 1 & 0 & 0 &  & w'x'yz'\\[0pt]
3 & 0 & 0 & 1 & 1 & 0 &  & w'x'yz\\[0pt]
4 & 0 & 1 & 0 & 0 & 0 &  & w'xy'z'\\[0pt]
5 & 0 & 1 & 0 & 1 & 1 & w+x'+y+z' & \\[0pt]
6 & 0 & 1 & 1 & 0 & 0 &  & w'xyz'\\[0pt]
7 & 0 & 1 & 1 & 1 & 1 & w+x'+y'+z' & \\[0pt]
8 & 1 & 0 & 0 & 0 & 1 & w'+x+y+z & \\[0pt]
9 & 1 & 0 & 0 & 1 & 1 & w'+x+y+z' & \\[0pt]
10 & 1 & 0 & 1 & 0 & 0 & w'+x+y'+z & \\[0pt]
11 & 1 & 0 & 1 & 1 & 1 &  & wx'yz\\[0pt]
12 & 1 & 1 & 0 & 0 & 0 & w'+x'+y+z & \\[0pt]
13 & 1 & 1 & 0 & 1 & 0 & w'+x'+y+z' & \\[0pt]
14 & 1 & 1 & 1 & 0 & 0 & w'+x'+y'+z & \\[0pt]
15 & 1 & 1 & 1 & 1 & 1 &  & wxyz\\[0pt]
\end{tabular}
\end{center}
\end{answer}
\end{subquestion}
\begin{subquestion}
find a minimal \href{boolean_algebra.org}{switching expression} for the function using a 4-variable \href{boolean_algebra.org}{karnaugh map}\\[0pt]
\begin{answer}
\begin{center}
\includesvg{/Users/mahmooz/brain/out/eUyRLTy}
\end{center}

so \(f(w,x,y,z)=xy'z+zxw+yz'w'+x'yz'\)\\[0pt]
\end{answer}
\end{subquestion}
\begin{subquestion}
find a minimal switching expression using a 3-variable karnaugh such that the values of w,z are the column variables and x is the only row variable\\[0pt]
\end{subquestion}
\end{question}
\begin{question}
given the function \(F\) that is described using the following karnaugh map\\[0pt]

\begin{center}
\includesvg{/Users/mahmooz/brain/out/6i2aZX2}
\end{center}
\begin{subquestion}
find \(f(A,B,C,D,E)\)\\[0pt]
\begin{answer}
step 1: \(A=0,\ B=0,\ B'=0,\ C=0\)\\[0pt]

\begin{center}
\includesvg{/Users/mahmooz/brain/out/bTs2oXP}
\end{center}

no possible implicants\\[0pt]
step 2: \(A=1\), the rest is \(0\)\\[0pt]

\begin{center}
\includesvg{/Users/mahmooz/brain/out/xexeHK3}
\end{center}

we get: \(ED\)\\[0pt]
step 3: \(B = 1\)\\[0pt]

\begin{center}
\includesvg{/Users/mahmooz/brain/out/BY1YakU}
\end{center}

we get: \(E'C'D' + E'CD\)\\[0pt]
step 4: \(B' = 1\)\\[0pt]

\begin{center}
\includesvg{/Users/mahmooz/brain/out/hhZ7kHK}
\end{center}

we get: \(ED\)\\[0pt]
step 5: \(A' = 1\)\\[0pt]

\begin{center}
\includesvg{/Users/mahmooz/brain/out/KgVxQ33}
\end{center}

we get: \(E'D\)\\[0pt]
if we summarize all the expressions we get we would arrive at the destination which is a switching expression describing the function \(f\)\\[0pt]

\[
  f(A,B,C,D,E) = ED + E'C'D' + E'CD + E'D
\]\\[0pt]
\end{answer}
\end{subquestion}
\end{question}
\begin{question}
given the following function\\[0pt]
\[
  f(A,B,C,D,E) = \sum(1,2,5,6,7,9,11,13,16,18,26,29,30) + {\sum}_\phi(19,20,22,23,25)
\]\\[0pt]
\begin{subquestion}
put the function in a karnaugh map so that each cell contains a minimal combination of the variables/constants \(\{A,B,0,1,\phi\}\)\\[0pt]
\begin{answer}
\begin{center}
\includesvg{/Users/mahmooz/brain/out/yYu8EYi}
\end{center}

first we build the original karnaugh map that fits this problem then we reduce it down to a smaller map with variables instead of constants 1 and 0\\[0pt]

\begin{tabular}{c|cccccccc|c}
  AB \textbackslash CDE & 000 & 001 & 011 & 010 & 110 & 111 & 101 & 100\\
  \hline
  00 & 0 & 1 & 3 & 2 & 6 & 7 & 5 & 4 & 0\\
  01 & 8 & 9 & 11& 10& 14& 15& 13& 12& 1\\
  11 & 24& 25& 27& 26& 30& 31& 29& 28& 3\\
  10 & 16& 17& 19& 18& 22& 23& 21& 20& 2\\
\end{tabular}

\begin{tabular}{c|cccccccc|c}
  AB \textbackslash CDE & 000 & 001 & 011 & 010 & 110 & 111 & 101 & 100\\
  \hline
  00 & 0 & 1    & 0    & 1 & 1    & 1    & 1 & 0    & 0\\
  01 & 0 & 1    & 1    & 0 & 0    & 0    & 1 & 0    & 1\\
  11 & 0 & $\varnothing$ & 0    & 1 & 1    & 0    & 1 & 0    & 3\\
  10 & 1 & 0    & $\varnothing$ & 1 & $\varnothing$ & $\varnothing$ & 0 & $\varnothing$ & 2
\end{tabular}

\(CDE = 000\)\\[0pt]

\begin{center}
\includesvg{/Users/mahmooz/brain/out/2LBxgrF}
\end{center}

we get: \(AB'\)\\[0pt]
\(CDE = 001\)\\[0pt]

\begin{center}
\includesvg{/Users/mahmooz/brain/out/PuZlq5P}
\end{center}

we get: \(A'\)\\[0pt]
\(CDE = 011\)\\[0pt]

\begin{center}
\includesvg{/Users/mahmooz/brain/out/gT9F99K}
\end{center}

we get: \(A'B\)\\[0pt]
\(CDE = 010\)\\[0pt]

\begin{center}
\includesvg{/Users/mahmooz/brain/out/avdXCnL}
\end{center}

we get: \(B' + A\)\\[0pt]
\(CDE = 110\)\\[0pt]

\begin{center}
\includesvg{/Users/mahmooz/brain/out/csgk3WS}
\end{center}

we get: \(B' + A\)\\[0pt]
\(CDE = 111\)\\[0pt]

\begin{center}
\includesvg{/Users/mahmooz/brain/out/MHlcB3z}
\end{center}

we get: \(B'\)\\[0pt]
\(CDE = 101\)\\[0pt]

\begin{center}
\includesvg{/Users/mahmooz/brain/out/ZYpOEK6}
\end{center}

we get: \(A' + B\)\\[0pt]
\(CDE = 100\) gives us \(0\) because the column is all \(0\)'s (that one \(\varnothing\) doesnt matter)\\[0pt]

\begin{center}
\includesvg{/Users/mahmooz/brain/out/J3OlntQ}
\end{center}
\end{answer}
\end{subquestion}
\end{question}
\end{document}